\section{Dispositif expérimental}

\subsection{Présentation du site expérimental}
L'Université du Sine Saloum El Hadj Ibrahima Niass (USSEIN)  dispose de plusieurs campus, dont son siège principal situé à Kaolack, dans le quartier de Sing Sing, à proximité directe de l'école franco-arabe sur la route de Sing Sing. C'est dans ce cadre que le champ d'application Agro TIC a été mis en place vers 2020--2021, en tant qu'espace expérimental dédié aux activités d'enseignement, de recherche et de démonstration en agriculture numérique ; il se situe au Sud-Est de l'université, et c'est à l'intérieur de ce champ, au fond, que se trouve notre parcelle expérimentale, dont le schéma ci-dessous présente une représentation. 
\begin{figure}[H]
    \centering
    \imgstd{images/03_site_experimental.jpg}
    \caption{Représentation cartographique du site expérimental. }
    \label{fig:site2}
\end{figure}

\subsection{Représentation de la surface expérimentale}

Notre surface d'expérimentation est représentée par le schéma ci-dessous. Elle est composée de 6 plots, répartis en 3 plots pour la zone test et 3 plots pour la zone témoin, chaque zone comportant un plot de gombo, un plot de navet et un plot de laitue. Les espacements entre les cultures varient selon les espèces : 30 cm pour le gombo, et 15 cm aussi bien pour le navet que pour la laitue. Pour la collecte des données, un capteur 8-en-1, fournissant huit paramètres du sol, est branché manuellement sur chaque plot test trois fois par jour, tandis qu'un capteur DHT22 vient compléter ce dispositif en mesurant la température et l'humidité ambiante du milieu.
\begin{figure}[H]
    \centering
    \imgstd{images/04_surface_experimentation.jpg}
    \caption{Représentation de la surface d'expérimentation. }
    \label{fig:surface2}
\end{figure}

\subsection{Méthode d’irrigation et d’apport en nutriment} 
\textbf{Zone Test :}

Dans la zone test, l'irrigation et l'apport en nutriments reposent sur une mesure des besoins réels de chaque culture. Trois plots y sont aménagés : un pour la laitue, un pour le navet et un pour le gombo chacun équipé d'un capteur mesurant huit paramètres du sol (humidité, température, teneur en sel, conductivité, pH, N, P et K), associé a l'utilisation de AccuWeather. Ce capteur, connecté à la machine, permet des relevés manuels deux fois par jour, facilitant ainsi la lecture des données. Les informations recueillies sont intégrées à la méthode de Penman-Monteith \cite{allen_crop_1998-1} conformément à la méthode standardisée FAO-56 \cite[chap.~2]{allen_crop_1998-1} afin d'estimer les besoins en eau des plantations et de déterminer en temps réel une éventuelle insuffisance ou un surplus d'eau ou d'éléments nutritifs. Ce calcul sert ainsi d'aide à la décision : l'irrigation et l'apport en nutriments ne sont effectués manuellement qu'en cas de besoin réel de la plante, de façon mesurée. 

 \textbf{Zone témoin :}
 
Dans la zone témoin, l'irrigation suit une méthode traditionnelle fixe, indépendante des besoins réels des plantes. Trois plots de $2\,\mathrm{m^2}$ y sont également aménagés : un pour la laitue, un pour le navet et un pour le gombo. L'arrosage y est effectué manuellement, chaque jour, matin et soir, selon un calendrier fixe qui ne tient compte ni des caractéristiques du sol ni de celles du climat, ni des besoins spécifiques de chaque culture.



